\sgaStableTarget{sga6.e05.continuation-of-2-9}We still suppose that @@M0@@. Multiplication by @@M2@@ under the law @@M1@@ defines a morphism of functors: @@M3@@, characterized by its value on the element @@M4@@ of @@M5@@. We therefore have:
@@M6@@
\sgaStableTarget{sga6.e05.eq.2-9-3}@@M7@@
We deduce:
\sgaStableTarget{sga6.e05.eq.2-9-4}@@M8@@
Indeed, the two members are additive functors equal on @@M9@@. From (2.9.3) and 1.15, it follows that the @@M10@@-th universal polynomial defining multiplication by @@M11@@ is isobaric of weight @@M12@@; it will be denoted by @@M13@@.

Finally, note that if @@M14@@ is a @@M15@@-@@M16@@-algebra, i.e. a @@M17@@-ring endowed with a @@M18@@-homomorphism: @@M19@@, and if @@M20@@, we have:
\sgaStableTarget{sga6.e05.eq.2-9-5}@@M21@@
This follows from @@M22@@, and from (2.9.4).

\sgaStableTarget{sga6.e05.sec.3}\section*{3. The operations @@M23@@}

We still work in the category @@M24@@.

\sgaStableTarget{sga6.e05.para.3-1}3.1. Let @@M25@@ be a @@M26@@-ring. From the operations @@M27@@ of @@M28@@, we shall define new operations, denoted by @@M29@@. They are intended to introduce a notion of Chern classes under conditions that will be specified later. For the moment, let us give a ``justification'' of the following definitions: if @@M30@@ is a vector bundle of rank @@M31@@, the @@M32@@-th Chern class of @@M33@@ will have to be (whatever theory of Chern classes is considered) the ``alternating sum of the exterior powers
